\chapter{Laparoscopic inguinal hernia repair (TAPP)}
\index{Laparoscopic inguinal hernia repair (TAPP)}

\section*{Definition and Overview}
Laparoscopic transabdominal preperitoneal (TAPP) inguinal hernia repair is a minimally invasive surgical technique used to repair groin hernias. During a TAPP procedure, the surgeon accesses the hernia defect by entering the peritoneal cavity with a laparoscope and instruments. The peritoneum overlying the hernia site is incised to expose the preperitoneal space, allowing the surgeon to reduce the herniated tissues (such as bowel or omentum) back into the abdomen. A synthetic mesh is then positioned over the myopectineal orifice of Fruchaud --- the common site for direct, indirect, and femoral hernias --- to reinforce the weakened abdominal wall. Finally, the peritoneal flap is sutured or tacked closed to cover the mesh, preventing it from contacting the abdominal organs and reducing the risk of adhesions.

\section*{Epidemiology}
Groin hernias are highly prevalent, with a lifetime risk of approximately 27\% in men and 3\% in women. Inguinal hernias represent about 75\% of all abdominal wall hernias. In many countries and globally, laparoscopic inguinal hernia repairs, including the TAPP and TEP (Totally Extraperitoneal) approaches, are widely performed. TAPP is particularly favoured for bilateral hernias, recurrent hernias where the previous repair was open, and in female patients where femoral hernias are more common.

\section*{Aetiology and Pathophysiology}
Inguinal hernias develop due to a combination of congenital predisposing factors and acquired weaknesses in the musculofascial structures of the groin. The pathophysiology involves the protrusion of abdominal contents through a defect in the transversalis fascia. Causes and risk factors include:

\begin{itemize}
    \item \textbf{Patent processus vaginalis:} A congenital failure of the peritoneal extension to obliterate during fetal development, creating a pre-formed pathway for indirect inguinal hernias.
    \item \textbf{Connective tissue degradation:} Age-related alteration in collagen metabolism, particularly the ratio of type I to type III collagen, which reduces the tensile strength of the groin fascia.
    \item \textbf{Chronic intra-abdominal pressure:} Elevated pressure secondary to chronic obstructive pulmonary disease (COPD) causing chronic cough, chronic constipation, bladder outlet obstruction (e.g.\ benign prostatic hyperplasia), and heavy lifting.
    \item \textbf{Prior abdominal surgery:} Previous surgical incisions that weaken the abdominal wall musculature.
    \item \textbf{Smoking:} Nicotine damages connective tissue and impairs collagen synthesis, significantly increasing hernia risk and recurrence rates.
\end{itemize}

\section*{Clinical Features}
A patient presenting with an inguinal hernia typically describes a palpable lump in the groin area. Clinical features include:

\begin{itemize}
    \item \textbf{Groin bulge:} A visible and palpable swelling in the groin or scrotum that increases in size with standing, coughing, or straining, and disappears or reduces when lying flat.
    \item \textbf{Discomfort or dull ache:} A dragging sensation in the groin that worsens towards the end of the day or after prolonged standing.
    \item \textbf{Localised tenderness:} Mild pain upon palpation of the hernia site.
    \item \textbf{Signs of incarceration or strangulation:} Severe, constant pain, skin erythema over the bulge, and systemic signs such as nausea and vomiting, indicating a surgical emergency.
\end{itemize}

\section*{Differential Diagnosis}
It is essential to differentiate an inguinal hernia from other anatomical abnormalities or sources of groin pain. These include:

\begin{itemize}
    \item \textbf{Hydrocele:} A fluid collection within the tunica vaginalis of the scrotum that transilluminates on physical examination.
    \item \textbf{Femoral hernia:} Herniation through the femoral canal, situated below and lateral to the pubic tubercle, carrying a high risk of strangulation.
    \item \textbf{Inguinal lymphadenopathy:} Enlarge lymph nodes in the groin secondary to localised infection, systemic inflammation, or haematological malignancies.
    \item \textbf{Varicocele:} Abnormally dilated veins of the pampiniform plexus within the spermatic cord, classically presenting as a soft 'bag of worms' swelling.
    \item \textbf{Osteitis pubis or groin strain:} Musculoskeletal inflammation of the pubic symphysis or adductor tendons, common in athletes, presenting with pain but no palpable hernia defect.
\end{itemize}

\section*{When to Seek Medical Care}
Patients with a known groin hernia should seek prompt medical advice if the hernia becomes increasingly painful or difficult to reduce.

\textbf{Contact your local emergency services for:}
\begin{itemize}
    \item Sudden, severe, or worsening groin or lower abdominal pain.
    \item A groin bulge that is firm, exquisitely tender, and cannot be pushed back (reduced) into the abdomen.
    \item Associated symptoms of bowel obstruction, such as nausea, persistent vomiting, severe abdominal bloating, and the inability to pass gas or stool.
    \item Discolouration (redness or purple skin changes) over the hernia site.
\end{itemize}

\section*{Diagnosis and Investigations}
The diagnosis of an inguinal hernia is predominantly clinical, but imaging can assist in ambiguous cases. Investigations include:

\begin{itemize}
    \item \textbf{Clinical examination:} Physical assessment of both groins while the patient is standing and lying down, utilising manual palpation during coughing (Valsalva manoeuvre) to detect a cough impulse.
    \item \textbf{Groin ultrasound:} A dynamic, non-invasive imaging modality that can identify the hernia sac, determine its contents, and differentiate indirect from direct hernias.
    \item \textbf{Computed tomography (CT) scan:} Useful in complex, giant, or recurrent hernias to define anatomical boundaries and check for occult hernias or bowel complications.
    \item \textbf{Magnetic resonance imaging (MRI):} Highly sensitive for identifying sports hernias, occult hernias, or musculoskeletal causes of chronic groin pain when clinical examination is inconclusive.
\end{itemize}

\section*{Management}
Management depends on the severity of symptoms and patient fitness. While asymptomatic hernias can be monitored, symptomatic hernias warrant surgical repair. Options include:

\begin{itemize}
    \item \textbf{Laparoscopic TAPP repair:} Minimally invasive surgery performed under general anaesthesia, using mesh reinforcement to repair the posterior wall of the inguinal canal via a transabdominal preperitoneal route.
    \item \textbf{Multimodal analgesia:} Postoperative pain management incorporating paracetamol, NSAIDs, and occasional short-term oral opioids to facilitate early mobilisation.
    \item \textbf{Watchful waiting:} Regular clinical surveillance for asymptomatic or minimally symptomatic hernias in patients with significant surgical comorbidities.
    \item \textbf{Activity modification:} Avoiding strenuous physical activity and heavy lifting (greater than 5--10 kg) for approximately 4 to 6 weeks post-surgery to allow tissue integration into the mesh.
\end{itemize}

\section*{Complications}
Laparoscopic TAPP repair is associated with a low complication rate, but specific risks must be considered:

\begin{itemize}
    \item \textbf{Chronic groin pain (inguinodynia):} Persistent neuropathic or nociceptive pain lasting more than three months postoperative, occurring in approximately 2--5\% of patients.
    \item \textbf{Seroma or haematoma:} A fluid or blood collection in the space previously occupied by the hernia sac, which typically resolves spontaneously.
    \item \textbf{Urinary retention:} Inability to void urine postoperatively, requiring temporary catheterisation, more common in elderly males.
    \item \textbf{Visceral injury:} Rare, accidental damage to the urinary bladder, small bowel, or colon during trocar insertion or preperitoneal dissection.
    \item \textbf{Vascular injury:} Damage to the inferior epigastric vessels or the iliac vessels (located in the 'doom triangle' under the peritoneal flap).
    \item \textbf{Hernia recurrence:} Re-emergence of the hernia, occurring in less than 2\% of cases when performed by experienced surgeons.
\end{itemize}

\section*{Prognosis and Outcomes}
The clinical outcomes for TAPP repairs are excellent, offering reduced postoperative pain, shorter hospital stays (often performed as a day procedure), and a faster return to work and normal activities compared to open mesh repairs. Long-term patient satisfaction is high. Recurrence rates are low, particularly when mesh of adequate size is deployed to cover all potential hernia spaces in the groin.

\section*{Prevention and Self-Care}
While congenital hernias cannot be prevented, patients can implement lifestyle changes to reduce the risk of hernia development or postoperative recurrence:

\begin{itemize}
    \item \textbf{Prevent chronic constipation:} Consume a high-fibre diet rich in fruits, vegetables, and whole grains, and maintain adequate fluid intake.
    \item \textbf{Utilise safe lifting mechanics:} Bend at the knees, keep the core engaged, and avoid holding your breath or straining during heavy lifts.
    \item \textbf{Maintain a healthy weight:} Reduces chronic intra-abdominal pressure on the inguinal wall.
    \item \textbf{Treat chronic respiratory conditions:} Manage asthma or chronic bronchitis promptly to avoid persistent coughing fits.
    \item \textbf{Smoking cessation:} Crucial for enhancing tissue healing, reducing postoperative infections, and preventing chronic cough.
\end{itemize}

\section*{Sources and Further Reading}
The following organisations provide reputable details on laparoscopic hernia repair:

\begin{itemize}
    \item Healthdirect Australia. \textit{Consumer guide on inguinal hernia symptoms, diagnosis, and surgical options.} https://www.healthdirect.gov.au/
    \item Mayo Clinic. \textit{Patient resource on inguinal hernia causes, surgical repair, and recovery tips.} https://www.mayoclinic.org/
    \item NHS. \textit{Overview of inguinal hernia repair, including laparoscopic and open surgical approaches.} https://www.nhs.uk/
    \item Royal Australasian College of Surgeons. \textit{Patient education pamphlets on groin hernias and laparoscopic surgery.} https://www.surgeons.org/
    \item British Hernia Society. \textit{Clinical consensus and patient guidelines regarding groin hernia management.} https://www.britishherniasociety.org/
\end{itemize}
